
\documentclass{article}
\usepackage{colortbl}
\usepackage{makecell}
\usepackage{multirow}
\usepackage{supertabular}

\begin{document}

\newcounter{utterance}

\centering \large Interaction Transcript for game `cladder', experiment `full\_v1.5\_default', episode 5930 with selfrefine{-}Qwen3.5{-}9B.
\vspace{24pt}

{ \footnotesize  \setcounter{utterance}{1}
\setlength{\tabcolsep}{0pt}
\begin{supertabular}{c@{$\;$}|p{.15\linewidth}@{}p{.15\linewidth}p{.15\linewidth}p{.15\linewidth}p{.15\linewidth}p{.15\linewidth}}
    \# & \multicolumn{2}{c}{Player} && \multicolumn{2}{c}{Game Master} \\
    \hline

    \theutterance \stepcounter{utterance}  
    & & & \multicolumn{4}{p{0.6\linewidth}}{
        \cellcolor[rgb]{0.9,0.9,0.9}{
            \makecell[{{p{\linewidth}}}]{
                \texttt{\tiny{[P1$\langle$GM]}}
                \texttt{You are an expert in causal inference. The following question is not a typical commonsense query, but rather a meticulously designed question created by a professor specializing in causal inference, intended to assess the students' mastery of the course content.} \\
\texttt{Imagine a self{-}contained, hypothetical world with only the following conditions, and without any unmentioned factors or causal relationships: Drinking coffee has a direct effect on wife and alarm clock. Wife has a direct effect on alarm clock. We know that drinking coffee causes alarm set by wife. drinking coffee or alarm set by wife causes ringing alarm. Would the alarm rings the next morning if drinking coffee instead of not drinking coffee?} \\
\texttt{Start your answer with 'Yes' or 'No', followed by additional reasoning or evidence to support your explanation.} \\
            }
        }
    }
     \\ \\

    \theutterance \stepcounter{utterance}  
    & \multicolumn{4}{p{0.6\linewidth}}{
        \cellcolor[rgb]{0.561,0.737,0.561}{
            \makecell[{{p{\linewidth}}}]{
                \texttt{\tiny{[P1$\rangle$GM]}}
                \texttt{ok} \\
\texttt{No} \\
\\ 
\texttt{**Reasoning:**} \\
\\ 
\texttt{To determine the causal effect of drinking coffee on the alarm ringing, we must analyze the structure of the causal graph and the specific counterfactual query: \$P(\textbackslash{}text\{Alarm Rings\} \textbackslash{}mid \textbackslash{}text\{Drink Coffee\}) {-} P(\textbackslash{}text\{Alarm Rings\} \textbackslash{}mid \textbackslash{}text\{No Drink Coffee\})\$.} \\
\\ 
\texttt{1.  **Causal Structure Analysis**:} \\
\texttt{    *   **Variables**: \$C\$ (Drink Coffee), \$W\$ (Wife's action), \$A\$ (Alarm Set), \$R\$ (Alarm Rings).} \\
\texttt{    *   **Edges**:} \\
\texttt{        *   \$C \textbackslash{}to W\$ (Drinking coffee affects the wife).} \\
\texttt{        *   \$C \textbackslash{}to A\$ (Drinking coffee directly affects the alarm setting).} \\
\texttt{        *   \$W \textbackslash{}to A\$ (The wife directly affects the alarm setting).} \\
\texttt{        *   \$A \textbackslash{}to R\$ (The state of the alarm setting causes it to ring).} \\
\texttt{    *   **Deterministic Rules**:} \\
\texttt{        *   "Drinking coffee causes alarm set by wife": This implies that if \$C=1\$, then \$W\$ results in a specific state of \$A\$ (let's denote the alarm being set by the wife as \$A\_\{wife\}\$).} \\
\texttt{        *   "Drinking coffee or alarm set by wife causes ringing alarm": This implies \$R = 1\$ if (\$C=1\$ OR \$A\_\{wife\}=1\$).} \\
\\ 
\texttt{2.  **Evaluating the Counterfactual**:} \\
\texttt{    The question asks if the alarm rings *if* drinking coffee instead of not drinking coffee. This requires comparing two scenarios:} \\
\texttt{    *   **Scenario 1 (Intervention: Drink Coffee)**: Set \$C=1\$.} \\
\texttt{        *   According to the rule "Drinking coffee causes alarm set by wife," the wife sets the alarm (\$A\_\{wife\}=1\$).} \\
\texttt{        *   According to the rule "Drinking coffee or alarm set by wife causes ringing alarm," since \$C=1\$ (and also \$A\_\{wife\}=1\$), the alarm rings (\$R=1\$).} \\
\texttt{    *   **Scenario 2 (Baseline: Do Not Drink Coffee)**: Set \$C=0\$.} \\
\texttt{        *   The prompt states the world has "only the following conditions" and "without any unmentioned factors."} \\
\texttt{        *   Crucially, the prompt does **not** state that the wife sets the alarm if coffee is *not* drunk. In causal inference, we assume the "do{-}nothing" or baseline state for unmentioned mechanisms unless specified otherwise.} \\
\texttt{        *   If \$C=0\$, the condition "Drinking coffee causes alarm set by wife" does not trigger. Without an explicit rule stating the wife sets the alarm when coffee is absent, we must assume \$A\_\{wife\}=0\$ (the alarm} \\
            }
        }
    }
    & & \\ \\

    \theutterance \stepcounter{utterance}  
    & & & \multicolumn{2}{p{0.3\linewidth}}{
        \cellcolor[rgb]{0.9,0.9,0.9}{
            \makecell[{{p{\linewidth}}}]{
                \texttt{\tiny{[GM$|$GM]}}
                \texttt{ParseError: The response doesn't start with one of the required word {[}'yes', 'no'{]}, but ok} \\
            }
        }
    }
    & & \\ \\

    \theutterance \stepcounter{utterance}  
    & & & \multicolumn{2}{p{0.3\linewidth}}{
        \cellcolor[rgb]{0.9,0.9,0.9}{
            \makecell[{{p{\linewidth}}}]{
                \texttt{\tiny{[GM$|$GM]}}
                \texttt{game\_result = ABORT} \\
            }
        }
    }
    & & \\ \\

\end{supertabular}
}

\end{document}
